\clearpage
\section{Box permutation through copying and contextual opening}
\label{sec:copying-box-permutation-overlaps}
This milestone proves the local copying overlap for BPerm and strengthens
its beta overlap to an exact comparison of whole-host execution costs.
Together with Section~\ref{sec:copying-box-permutation-opening}, the two local
compute diagrams now cover the shallow lambda cases identified in the paper's
\href{https://arxiv.org/html/1907.01257v5\#A8.SS3.SSS6}{Appendix~H.3.6}:
a contraction tree copying a lambda payload, and application opening its box.
The copying rule is still Figure~17, and the application rule retains its
stable-argument premise. The full operation alphabet, original program terms
and Figure~16 translation are unchanged.

There is an essential difference between the endpoints of the two diagrams.
Beta opening eliminates the compared outer boxes and gives recursively
isomorphic successor states. Copying produces two surviving BPerm pairs.
The new proof assembles those pairs into a structural residual derivation and
tracks their distinct actual edges. It does not identify that derivation with
observational equivalence, or with a fully certified residual specimen.
At this stage the recursive residual specimen certificates were still open.
Section~\ref{sec:copying-residual-specimens} now supplies them for copying,
including deep and mixed certificate assembly. Exhaustive overlap transfer
and universal forward/converse BPerm robustness remain unfinished. This
section's structural diagram alone does not prove Proposition~12.2.

\Needspace{14\baselineskip}
\subsection{The complete copy reduct and its boundary action}
\termdef{copying-reduct-boundary-naturality}
Write $E_{k,h}$ for the ordered environment interface of $k$ term wires followed
by $h$ atom wires. Let $T$ be any genuine term contraction tree with $n+1$
inputs, and let $G,H:*\to E_{k,h}$ be payload-shaped graphs with a typed boundary
isomorphism $f:G\to H$. A boundary isomorphism preserves ordered incidences,
wire types and boundary roles but may permute external positions. Its singleton
input action is necessarily the identity; write $p=f_{\mathrm{out}}$ for its
free-output action.

The literal reduct $A_T(G)$ contains an active copy of $G$, a retained copy fed
by the residual tree, and the original two-use variable and atom distributors.
The residual tree is the matched tree $T$ with its selected input supplied by
a weakening. It is not replaced by a canonical tree. Naturality constructs a
boundary isomorphism $A_T(G)\to A_T(H)$ fixing all $n+1$ inputs and acting by
$p$ on the final outputs. This statement concerns graph assembly; it does not
assert that arbitrary payloads or their external permutations are copyable.

The auxiliary \texttt{reboundary} operation reindexes equal boundary lists
without changing the edge or internal carriers. Its isomorphism uses the
original port and edge maps after equality elimination; source and target
preservation reduce to the identity boundary cast. The three transport lemmas
state explicitly that the edge bijection is unchanged and input/output actions
are the original actions with the corresponding finite-index casts.

\lean{Copying/CopyBoundaryNaturality}{8}{19}{BoundaryIso.reboundary}
\lean{Copying/CopyBoundaryNaturality}{21}{23}{BoundaryIso.reboundary_edges}
\lean{Copying/CopyBoundaryNaturality}{25}{29}{BoundaryIso.reboundary_input}
\lean{Copying/CopyBoundaryNaturality}{31}{35}{BoundaryIso.reboundary_output}

\Needspace{12\baselineskip}
The retained branch is $R_T;G$, where $R_T$ is the original residual tree.
Glue its identity boundary isomorphism to $f$. Compatibility at the common
singleton interface follows because $f$ fixes the sole input. The resulting
\texttt{retained} map fixes every residual input and inherits $f$'s output action.
The \texttt{raw} construction then assembles $f$ on the active branch and this
retained map on the other branch. It uses the original Boolean branch family
and the original middle-block output shuffle. Finally \texttt{copies} reindexes
the equal input and output lists into the interfaces used by Figure~17.

\lean{Copying/CopyBoundaryNaturality}{44}{45}{CopyRule.Boundary.retained}
\lean{Copying/CopyBoundaryNaturality}{47}{49}{CopyRule.Boundary.retained_input}
\lean{Copying/CopyBoundaryNaturality}{51}{53}{CopyRule.Boundary.retained_output}
\lean{Copying/CopyBoundaryNaturality}{55}{61}{CopyRule.Boundary.raw}
\lean{Copying/CopyBoundaryNaturality}{63}{64}{CopyRule.Boundary.copies}

\Needspace{12\baselineskip}
For a branch $b$ and an original output index $j$, let $\iota_b(j)$ be its
position after that shuffle. The output lemma proves
\[
 \iota_b(j)\longmapsto\iota_b(p(j)).
\]
It follows by the parallel embedding equation, using $f$ on the active branch
and the retained-output lemma on the other. The input lemma splits each input
index into its Boolean branch and local position. In the active branch it uses
the singleton-input identity; in the retained branch it uses the residual-input
identity. Joining the split index restores the original position. This proves
pointwise preservation of the active input and every unselected input, including
$n=0$.

\lean{Copying/CopyBoundaryNaturality}{67}{86}{CopyRule.Boundary.copies_output}
\lean{Copying/CopyBoundaryNaturality}{89}{115}{CopyRule.Boundary.copies_input}

The existing typed distributor automorphism acts on each of its two uses by
$p$ and on its final interface by $p$. Every incoming index is some
$\iota_b(j)$, so the preceding output equation proves the compatibility needed
to glue the two-copy graph to that distributor. The definition \texttt{after}
is precisely this serial gluing. Its outer input action is inherited from the
copies, and its outer output action from the distributor.

\lean{Copying/CopyBoundaryNaturality}{118}{126}{CopyRule.Boundary.compatible}
\lean{Copying/CopyBoundaryNaturality}{129}{130}{CopyRule.Boundary.after}
\lean{Copying/CopyBoundaryNaturality}{132}{136}{CopyRule.Boundary.after_input}
\lean{Copying/CopyBoundaryNaturality}{138}{142}{CopyRule.Boundary.after_output}

\Needspace{16\baselineskip}
The edge laws give additional information beyond existence of an isomorphism.
An edge in the active copy maps to the corresponding active-copy edge under
$f$; an edge in the retained copy maps to the retained-copy edge under the same
$f$. Every residual-tree and weakening edge is fixed. These follow by unfolding
the serial and parallel edge maps and the proved invariance of the edge map
under boundary reindexing. Hence the map neither interchanges the two copies
nor changes the matched tree presentation.

\lean{Copying/CopyBoundaryNaturality}{145}{147}{CopyRule.Boundary.after_leftEdge}
\lean{Copying/CopyBoundaryNaturality}{150}{152}{CopyRule.Boundary.after_rightEdge}
\lean{Copying/CopyBoundaryNaturality}{155}{161}{CopyRule.Boundary.after_residualEdge}

For a typed external permutation $p$, specialise $f$ to the canonical boundary
isomorphism $G\to pG$. Realise the resulting output action by an external
permutation and invert the comparison. This gives a boundary-fixing recursive
isomorphism
\[
 A_T(pG)\cong pA_T(G).
\]
All original distributors remain in both graphs. In particular, this is not
an additional copying rule with a weakened copyability premise.

\lean{Copying/CopyBoundaryNaturality}{165}{177}{CopyRule.Boundary.permuted}

\Needspace{14\baselineskip}
\subsection{A literal BPerm pair through two copies}
\termdef{copying-two-copy-derivation}
Recall that \texttt{BoxPermutation.Derivation} is the explicit structural
closure of literal BPerm steps and recursive isomorphisms. Its constructors
alone make no semantic soundness claim. The new boundary-reindexing lemma
follows by equality elimination. Its \texttt{copies} lemma applies the same
supplied derivation once to the active payload and once to the payload behind
the unchanged residual tree. Parallel assembly and boundary reindexing finish
that proof. The \texttt{copyingAfter} lemma appends the identity derivation of
the unchanged distributor. These proofs expose both uses of the supplied pair.

\lean{Copying/BoxPermutationCopying}{12}{14}{BoxPermutation.Derivation.reboundary}
\lean{Copying/BoxPermutationCopying}{18}{25}{BoxPermutation.Derivation.copies}
\lean{Copying/BoxPermutationCopying}{27}{29}{BoxPermutation.Derivation.copyingAfter}

Let $B,C:*\to E_{k+1,h}$ be bodies related by a boundary isomorphism $f$ which
fixes the bound output and acts as $p$ on the free outputs. The condition is
exactly \texttt{Opening.Action} from the preceding section. Let $L(B)$ denote
the original lambda head followed by its original box. The compensated payload
is the lambda head followed by $p^{-1}$ applied to the free outputs of $C$'s
box. Moving this compensation past the lambda head gives the recursive
isomorphism from the compensated payload to $p^{-1}L(C)$.

\Needspace{20\baselineskip}
The body action supplies a genuine BPerm pair between the two boxes. Prepending
the common lambda head gives \texttt{payload\_pair}; prepending the common
contraction tree and moving the compensation outside gives
\texttt{source\_pair}. Applying the two-copy closure lemma gives a residual
derivation to the reduct of the compensated payload. Copy-construction
isomorphism invariance, followed by the new boundary naturality theorem, moves
the compensation outside the whole reduct. This is \texttt{residual\_pair}.

\lean{Copying/BoxPermutationCopying}{38}{42}{BoxPermutation.Copying.compensatedPayload}
\lean{Copying/BoxPermutationCopying}{44}{45}{BoxPermutation.Copying.payloadIso}
\lean{Copying/BoxPermutationCopying}{47}{49}{BoxPermutation.Copying.payload_pair}
\lean{Copying/BoxPermutationCopying}{51}{55}{BoxPermutation.Copying.source_pair}
\lean{Copying/BoxPermutationCopying}{59}{64}{BoxPermutation.Copying.residual_pair}

\Needspace{14\baselineskip}
\subsection{The actual copying transitions and surviving boxes}
\termdef{copying-bperm-copy-diamond}
Both $L(B)$ and $L(C)$ satisfy the original copyability predicate by its lambda
constructor, independently of their bodies. Consequently the diagram has actual
compute transitions
\[
 \begin{aligned}
  \operatorname{before}_T(L(B))&\longrightarrow A_T(L(B)),\\
  p^{-1}\operatorname{before}_T(L(C))&\longrightarrow p^{-1}A_T(L(C)).
 \end{aligned}
\]
Both tokens start in rewrite mode at the selected tree input and finish there
in search mode. The first transition is the original local copying rule;
the second is that same rule transported through the already-proved external
permutation context. The theorem also supplies the structural residual
derivation just constructed. It never applies the copyability predicate to an
arbitrarily externally permuted operation corolla.

Prepending the genuine contraction-search transition gives
\texttt{initial\_diamond}: exactly two original machine transitions on each
side, beginning in search mode. No termination or stability condition on the
box body is needed, and $T$ is an arbitrary contraction-tree presentation.

\lean{Copying/BoxPermutationCopying}{68}{79}{BoxPermutation.Copying.copy_diamond}
\lean{Copying/BoxPermutationCopying}{83}{97}{BoxPermutation.Copying.initial_diamond}

The \texttt{boxEdge} coordinate is the box's original corolla edge inserted to
the right of the lambda head. The residual-label theorem instantiates the
literal copying label law at this edge: both copies retain the complete original
box label, including nested contents. The disjointness theorem instantiates the
copy construction's distinct-branch theorem, so the two occurrences remain
different even when their labels are equal.

Finally, choose $C$ to be the original body $B$ with the canonical bound-fixing
extension of $p$. Its action condition holds by construction. The theorem
\texttt{permuted\_copying} therefore gives both source and residual structural
derivations without a further body-action hypothesis.

\lean{Copying/BoxPermutationCopying}{100}{101}{BoxPermutation.Copying.boxEdge}
\lean{Copying/BoxPermutationCopying}{105}{109}{BoxPermutation.Copying.residual_box_labels}
\lean{Copying/BoxPermutationCopying}{111}{113}{BoxPermutation.Copying.residual_boxes_distinct}
\lean{Copying/BoxPermutationCopying}{116}{124}{BoxPermutation.Copying.permuted_copying}

\Needspace{14\baselineskip}
\subsection{Cancelling equal prefixes in whole-host executions}
\termdef{copying-equal-prefix-costs}
Write $\mathsf{Term}_c(s)$ for termination of state $s$ in exactly $c$ original
machine transitions. Proved successor coherence gives, for every actual prefix
$s\xrightarrow{n}u$,
\[
 \mathsf{Term}_c(s)\quad\Longleftrightarrow\quad
 \exists r,\ c=r+n\ \land\ \mathsf{Term}_r(u).
\]
This is a property of the concrete machine, not an added determinism premise.
If also $t\xrightarrow{n}v$ and $u,v$ are recursively isomorphic as states,
transport of exact successful costs across that isomorphism identifies the
right-hand expressions. Thus the sources have the same exact successful costs.
The theorem \texttt{equal\_prefix\_costs} is this argument. To obtain its
\texttt{plug} version, first lift both actual prefixes and their endpoint state
isomorphism through the typed host slot, then apply the argument to the whole
host machine. This order avoids treating local return as global finality.

\Needspace{16\baselineskip}
The state interface-cast operation, its run lemma and its isomorphism lemma
only eliminate equalities between ordered boundary lists. They make the slot
interface equations explicit without inserting transitions or modifying graphs.

\lean{Copying/EqualPrefixCosts}{11}{12}{Combined.castInterfaces}
\lean{Copying/EqualPrefixCosts}{14}{17}{Combined.run_castInterfaces}
\lean{Copying/EqualPrefixCosts}{19}{22}{Combined.StateIso.castInterfaces}
\lean{Copying/EqualPrefixCosts}{26}{30}{Combined.equal_prefix_costs}
\lean{Copying/EqualPrefixCosts}{35}{41}{Combined.equal_prefix_costs_plug}

\Needspace{12\baselineskip}
A stronger cancellation statement does not require endpoint isomorphism:
\[
 \bigl(\forall c,\ \mathsf{Term}_c(s)\leftrightarrow\mathsf{Term}_c(t)\bigr)
 \ \Longleftrightarrow\
 \bigl(\forall r,\ \mathsf{Term}_r(u)\leftrightarrow\mathsf{Term}_r(v)\bigr).
\]
The reverse implication applies the prefix equation on both sides. For the
forward implication, inspect the source comparison at cost $r+n$ and cancel
$n$ in the resulting natural-number equalities. The two \texttt{at} variants
combine these arguments with explicit interface casts and actual plugging.
They apply to arbitrary typed outer slots, without requiring local finality to
be preserved by insertion.

\lean{Copying/EqualPrefixCosts}{45}{56}{Combined.equal_prefix_spectra_iff}
\lean{Copying/EqualPrefixCosts}{59}{68}{Combined.equal_prefix_costs_at}
\lean{Copying/EqualPrefixCosts}{72}{82}{Combined.equal_prefix_spectra_at}

\Needspace{14\baselineskip}
\subsection{Quantitative conclusions for the two overlaps}
\termdef{copying-contextual-redex-costs}
For beta opening, the compensated graph isomorphism from the preceding section
also preserves the search mode and the selected input token; this defines
\texttt{successorIso}. Each source reaches its successor in six transitions.
Equal-prefix cancellation therefore proves \texttt{redex\_costs}, equality of
all exact successful costs of the two original application sources. Its
\texttt{at} version places the entire comparison in an arbitrary matching host
slot; existentially quantifying the common cost gives termination equivalence.
The stable argument is still a premise of both actual beta prefixes. The body
may contain arbitrary operations or nested boxes, and need not terminate.

These hosted statements start with the token at the inserted application's
input. They are stronger than equality of isolated successor costs, but do not
claim that every evaluation of an arbitrary box context has already been
classified into this local case. That classification remains part of template
robustness.

\lean{Copying/BoxPermutationOverlapCosts}{15}{21}{BoxPermutation.Opening.successorIso}
\lean{Copying/BoxPermutationOverlapCosts}{25}{30}{BoxPermutation.Opening.redex_costs}
\lean{Copying/BoxPermutationOverlapCosts}{35}{45}{BoxPermutation.Opening.redex_costs_at}
\lean{Copying/BoxPermutationOverlapCosts}{49}{58}{BoxPermutation.Opening.redex_termination_at}

For copying, the two-transition prefixes give a different, precisely delimited
conclusion. In every matching host, equality of all source costs is equivalent
to equality of all residual costs. This is \texttt{residual\_spectra\_at}.
It retains the semantic residual comparison as an obligation; the structural
derivation proved above does not discharge it by itself. This is the cost
handoff needed when the residual pairs are eventually placed into the general
robustness argument.

\lean{Copying/BoxPermutationOverlapCosts}{69}{84}{BoxPermutation.Copying.residual_spectra_at}

\Needspace{14\baselineskip}
\subsection{Concrete boundary and context checks}
The regression body from the preceding section contains a nested lambda,
an application and an atom use. Its permutation swaps two free variables and
two atoms while fixing the bound output. \texttt{changedBody} uses the actual
interface permutation of this original translation. The \texttt{nested\_copying}
theorem instantiates the two-transition diagram for every genuine contraction
tree, retaining its residual derivation. The binary-tree check proves the two
box occurrences distinct; the unary-tree check covers no residual external
inputs. A three-input tree checks that input position two is fixed by the
complete reduct's boundary action.

\lean{Copying/BoxPermutationOverlapExamples}{11}{11}{BoxPermutation.OverlapExamples.changedBody}
\lean{Copying/BoxPermutationOverlapExamples}{15}{23}{BoxPermutation.OverlapExamples.nested_copying}
\lean{Copying/BoxPermutationOverlapExamples}{26}{30}{BoxPermutation.OverlapExamples.binary_boxes}
\lean{Copying/BoxPermutationOverlapExamples}{33}{37}{BoxPermutation.OverlapExamples.unary_copying}
\lean{Copying/BoxPermutationOverlapExamples}{40}{44}{BoxPermutation.OverlapExamples.third_input_fixed}

\Needspace{18\baselineskip}
The full redex cost comparison is instantiated with the nested body and the
original stable unit argument. The host is a real unary-negation operation
followed by a typed hole, whose edge coordinate is supplied explicitly.
The hosted cost theorem follows from the general hosted beta comparison with
reflexive interface equalities. In this host the inserted application's input
is internal, so the check exercises the situation in which local return is
not global termination.

\lean{Copying/BoxPermutationOverlapExamples}{47}{51}{BoxPermutation.OverlapExamples.nested_redex_costs}
\lean{Copying/BoxPermutationOverlapExamples}{55}{57}{BoxPermutation.OverlapExamples.host}
\lean{Copying/BoxPermutationOverlapExamples}{59}{59}{BoxPermutation.OverlapExamples.hole}
\lean{Copying/BoxPermutationOverlapExamples}{61}{68}{BoxPermutation.OverlapExamples.hosted_redex_costs}

\Needspace{10\baselineskip}
Finally \texttt{selfApplication} is the scoped syntax
$(\lambda x.x\,x)(\lambda x.x\,x)$ in the same nonempty free environment.
The final regression applies the residual derivation theorem to its literal
translation and the same nontrivial free permutation. This checks that nested
self-application is retained and that no body-termination premise has entered
the copying proof. It is not a separate divergence theorem for that syntax.

\lean{Copying/BoxPermutationOverlapExamples}{72}{73}{BoxPermutation.OverlapExamples.selfApplication}
\lean{Copying/BoxPermutationOverlapExamples}{75}{81}{BoxPermutation.OverlapExamples.self_application_copying}
